% Short variant of sec_tiers_elicitation.tex (S5, S6). The opus_alias
% diagnostics, the ambition diagnostic, Figures fig:packings and fig:armT,
% Table tab:tracegrid, the P-T3 limits arithmetic, the bundling disclosure,
% the falsifier discussion and the faithfulness exclusion list are relocated
% verbatim to Appendix E; nothing is dropped.
\section{The tier ladder: three attractor families}

Holding prompt, container and scoring fixed, we varied only the nominal proposer tier across the
three cells that discriminate hardest between prediction and family optimum --- $N = 13$, 21 and
31, each inside a trap zone. The tiers do not merely differ in success rate: they attempt
qualitatively different constructions (Figure~\ref{fig:packings}, Appendix~E.1).

\textbf{The weak tier truncates templates.} The Haiku-tier proposer produces uniform grids of
radius $1/(2k)$ truncated to $N$ circles, with corner fillers only when the grid underfills. Across
the 45 original bare invocations spanning $N \in \{13, 17, 21, 31, 35, 37, 43\}$, \textbf{35 were
geometrically valid at the primary $10^{-6}$ tolerance (77.8\%), 29 at $10^{-9}$ (64.4\%)} ---
recomputed from the ledger, superseding the figure in earlier drafts --- and at the three
discriminating cells 12 of 16 valid samples landed on the predicted value, the rival being reached
twice, both at $N = 21$.

\textbf{The middle tier perturbs and mixes.} The Sonnet-tier arm was preregistered in
\texttt{arm\_s\_preregistration.txt}, disclosing that 5 of 20 samples at $N = 13/21$ had been seen
before registration and that $N = 31$ was fully blind. It was valid 30/30 at $10^{-6}$ and 27/30
(90\%) at $10^{-9}$ --- so the ``100\%'' leg is tolerance-dependent and both numbers are stated ---
and almost never on-prediction: 0/10, 0/10, 1/10, so 1/30 pooled against the weak tier's 12/16 at
the same cells. Instead of truncating a uniform grid it perturbs one, with enlarged edge rows,
hexagonal interior rows and two or three distinct radii in the same packing (29/30 multi-radius).
One Sonnet sample at $N = 31$ summed to 2.7499999991, above 2.7485281, the best the family reaches
there, and is the only sample in the swept corpus to leave the recipe family upward (arm M's
converge cells hold three more, \S3.4). That sample also exposes a scoring defect: the two values
differ by $1.47\times10^{-3}$, \emph{inside} the registered $2\times10^{-3}$ window, so the value
classifier counted the escape as a hit on the construction it exceeded. Classified by structure
instead, Sonnet's rival count at $N = 31$ is 2/10 and pooled 5/30, and \textbf{every categorical
claim in this paper about ``reached the rival'' or ``left the family'' is made at structure level
or at $10^{-6}$, never at $2\times10^{-3}$} (Appendix~E.2).

\textbf{The top tier attempts recursive constructions and fails to build them.} The third arm was
invoked through a bare tier alias. We name it \texttt{opus\_alias} throughout and make no claim
about which weights served it: the runtime accepts only the bare alias and exposes no dated
identifier, so the binding is a vendor promise rather than an attestable fact, and no version
number appears anywhere in this paper; completion-time and token-count anomalies were observed in
the unreleased session transcript, but the ledger carries no latency or token field, so nothing
here relies on them. Under those caveats the arm --- preregistered fully blind in
\texttt{arm\_o\_preregistration.txt} --- was valid in 4 of 30 invocations (13\%) at both tolerances,
attempting at every cell a more ambitious construction than either other tier: quarter-circle
corners at $r = 0.25$ with Apollonius-style fillers at $N = 13$ (10/10 samples), mixed-radius
$4 \times 4$-ish grids at $N = 21$, and a coarse $3 \times 3$ grid at $r \approx 1/6$ with border
strips at $N = 31$ (0/10 valid). Valid samples score \emph{below} the trap they were expected to
fall into (1.26--1.41 at $N = 13$ against 1.625), and a disclosed post-hoc diagnostic finds no
tolerance artifact behind the 13\% (Appendix~E.3).

Two decisions here were made after seeing the data and are recorded as deviations
(Table~\ref{tab:deviations}). P-O1, P-O2 and P-O4 are reported \textbf{not evaluable}, on the
grounds that a validity collapse of this size makes an on-prediction tier comparison
uninterpretable; P-O3 is trivially satisfied on 4/4 valid samples; the registered disconfirmation
--- regression toward the trap --- did not occur. De-registering three of four predictions in the
one arm that disconfirmed is exactly the move preregistration exists to constrain, which is why it
is tabled rather than narrated. The arm was also initially excluded from the ladder, then included
once its results were known; we follow the later decision and report the ladder both ways
(Table~\ref{tab:ladder}) --- with the arm 83\% $\rightarrow$ 100\% $\rightarrow$ 13\% on matched
cells (78\% over all seven weak-tier cells), without it 83\% $\rightarrow$ 100\%.

\begin{table}[t]
\centering
\caption{Three-attractor ladder. Valid columns count samples passing containment and non-overlap at the $10^{-6}$ and $10^{-9}$ tolerances; On-pred and Rival count valid samples landing on the predicted value and on the rival argmax.}
\label{tab:ladder}
\footnotesize
\setlength{\tabcolsep}{4pt}
\begin{tabular}{@{}p{1.7cm}p{1.2cm}cp{2.2cm}p{1.0cm}p{1.0cm}p{1.5cm}cp{2.0cm}@{}}
\toprule
Tier & Cells & $n$ & Attempted family & Valid $10^{-6}$ & Valid $10^{-9}$ & On-pred & Rival & Failure mode \\
\midrule
haiku (bare) & 13, 17, 21, 31, 35, 37, 43 & 45 & uniform grid, truncated; fillers when underfull & 35 (78\%) & 29 (64\%) & 12/16 (3 matched cells) & 2 & overlap 6, outside 1, parse 3 \\
haiku, matched cells & 13, 21, 31 & 60 & as above & 50 (83\%) & --- & 35/50 & 2 & overlap 5, outside 2, parse 3 \\
sonnet & 13, 21, 31 & 30 & perturbed grid, hexagonal rows, 2--3 radii & 30 (100\%) & 27 (90\%) & 1/30 & 5 (structure) & none \\
\texttt{opus\_alias} $\ddagger$ & 13, 21, 31 & 30 & quarter-circle corners, Apollonius fillers, coarse grid + border strips & 4 (13\%) & 4 (13\%) & 0/4 & 0 & overlap 24, nonpositive radius 2 \\
\bottomrule
\end{tabular}
\end{table}

$\ddagger$ Addressed through a bare serving alias; the alias-to-weights binding is not attestable,
and this row is not a statement about any released model version. Tier denominators are not matched
(the weak tier spans seven $N$ against three), which is why the matched-cell row is given
separately and why the matched comparison (83\% $\rightarrow$ 100\% $\rightarrow$ 13\%) is the one to
quote across tiers.

``Monotone'' applies to one axis only: ambition rises monotonically --- truncated template,
perturbed hybrid, recursive gasket --- while validity rises then collapses, and ``ambition'' is
measurable rather than a qualitative read, a disclosed post-hoc diagnostic being monotone on both
of its operationalizations (Appendix~E.4). At the canonical and plausibly contaminated cell
$N = 26$ all tiers converge on the same 2.5414 attractor; they diverge only at withheld trap cells.
The branch rule of \S2 is therefore a weak-tier regularity --- cross-vendor under two instruments
(GM3's pooled bar carried by its non-discriminating cells, \S3.6),
family-bounded in both directions, cell-level modes confirmed in one lineage (\S3.6) --- and this
is a \textbf{boundary condition on the main result}, not an independent second finding. It also
cuts against \S1's framing: the discovery systems cited there run frontier or mid-tier proposers,
so the tier at which our regularity holds is the tier they do not use, while the tier closest to
what they use escapes the trap in our own data.

\section{Elicitation and faithfulness}

This section is explicitly secondary material and the result leans negative: the registered
validity prediction fails, the concentration effect survives only as a fragile directional
finding, and the faithfulness audit checks description against artifact, not against process.

\subsection{Design and bundling disclosure}

A ten-versus-ten pilot at $N = 21$ compared the bare prompt against a variant asking the proposer
to name its construction on a leading \texttt{METHOD:} line: validity rose from 7/10 to 10/10 and
the higher-scoring rival disappeared, from 2 of 7 valid bare samples to 0 of 10 trace samples. Four
predictions and an explicit falsifier were registered before any scaled sample was drawn
(\texttt{arm\_t\_preregistration.txt}, sha256 \texttt{ab7900a8\ldots}), disclosing that the pilot's
prompt had drifted from the bare template beyond the method line, so the pilot's intervention was
method-line-plus-rewording, bundled; pilot samples are never pooled with \texttt{trace\_v2}. The
scaled prompt, \texttt{trace\_v2}, is itself a \textbf{bundled prompt-format-and-trace-request}
intervention whose output-line rewording is not cosmetic --- \S6.2 shows it doing measurable work
--- so the measured effect is \emph{method-line-plus-output-line-rewording}, a weaker version of
the confound for which the pilot was discarded, held to the same standard (Appendix~E.9). The
scaled arm ran 100 new invocations, bringing both arms to 20 per cell at $N \in \{13, 21, 31\}$;
the corpus arithmetic, from 215 logged square invocations to the study's 858 unconditioned-and-
elicitation invocations (678 before arms CP and RP, \S3.9), is in Appendix~E.5. One disclosure: \textbf{20 of the 60 bare samples are
pre-existing arm-F rows} whose results were known when P-T1--P-T3 were written, so only trace\_v2
is fully fresh (\S6.4).

\subsection{Registered outcomes}

Preregistered criteria are evaluated exactly as registered, and unregistered alpha thresholds are
applied nowhere. The paired per-cell grid is Table~\ref{tab:tracegrid} (Appendix~E.6).

\textbf{P-T1, validity: NOT confirmed.} P-T1 is the one prediction that registered an alpha
(``validity $\ge$ bare at EACH $N$, and pooled one-sided Fisher $p < 0.05$''). Direction held at all
three cells, but the pooled test gives $p = 0.30$: 53/60 (88\%) trace\_v2 against 50/60 (83\%) bare,
a difference whose detection needs several hundred samples per arm at 80\% power --- a failure to
detect, not a demonstration of no effect. The direction that did hold is not geometric: bare had
\textbf{3 parse failures and 7 geometric failures}, trace\_v2 \textbf{0 parse and 7 geometric}, so
among parsed samples geometric validity is 50/57 (87.7\%) versus 53/60 (88.3\%) --- the P-T1
direction is format compliance, exactly what the reworded output-format line manipulates.

\textbf{P-T2, rival suppression: CONFIRMED as registered} (directional, no alpha): 1 of 53 valid
trace\_v2 against 2 of 50 valid bare. Counts this near zero carry no inferential weight and we
claim none (Fisher $p = 0.48$). One trace\_v2 sample at $N = 31$ hit the rival 2.7485281 to within
$3.6\times10^{-8}$, the first weak-tier rival hit at $N = 31$ in any arm.

\textbf{P-T3, anchor concentration: met as registered (directional), inferentially fragile.} Among
valid samples, 46 of 53 trace\_v2 (87\%) landed on the registered prediction against 35 of 50 (70\%)
bare. The $p$-value it never carried is 0.0325, one-sided uncorrected (\S6.4).

\textbf{P-T4, faithfulness: confirmed under a corrected scorer.} The registered scorer gives 38 of
41 scoreable claims matching (93\%) against a 90\% threshold; blind hand-adjudication under a rubric
frozen before rescoring gives 54 of 56 (96.4\%). \S6.5 reports both.

\subsection{The registered falsifier, both readings}

The registered falsifier reads: \emph{``if trace\_v2 validity \texttt{<=} bare at 2+ of 3 N, the
pilot effect was the bundled rewording, not the method-line request --- reported as such, not
reframed.''} Observed validity: $N = 13$, 18/20 both arms --- \textbf{tie}; $N = 21$, trace\_v2
18/20 against bare 15/20; $N = 31$, 17/20 both arms --- \textbf{tie}. A tie satisfies \texttt{<=},
so under the registered wording the falsifier is \textbf{triggered at 2 of 3 $N$} and we report it
as triggered; the code instead used strict \texttt{<}, under which it fires at 0 of 3 cells --- a
convention that appears nowhere in the preregistration and read the same tie favorably twice. We
name the failure mode \textbf{operationalization drift between preregistration text and analysis
code}. Under the authoritative tie-inclusive reading, \textbf{the pilot's validity effect is
attributed to the bundled rewording, not to the method-line request}, and P-T1 fails under both
readings regardless (Appendix~E.10).

\subsection{Limits of P-T3: what the concentration result will and will not carry}

P-T3 met its registered directional criterion. It is not a flagship result, on four counts worked
through with their arithmetic in Appendix~E.7: it registered \emph{no alpha}, so the $p$-value is a
post-hoc addition; it fails \emph{multiplicity} correction, since Holm's first threshold is
$0.05/3 = 0.0167$ and $p = 0.0325$ also fails Bonferroni and Benjamini--Hochberg FDR; it is
\emph{concentrated in one cell}, the pooled 17-point gap being carried by $N = 13$ and driven by a
low bare rate there rather than a high trace rate, with $p = 0.31$ once that cell is dropped; and
it is \emph{wave-confounded}, the control arm's own between-wave drift being comparable to the
effect attributed to the manipulation. So: a minimal method-naming request does not make the
proposer better at geometry and does not measurably change which rare constructions it reaches; in
the registered direction it concentrates output onto the anchor, but that concentration fails
correction, is carried by one cell, is confounded with collection wave, and is bundled with an
output-format rewording that demonstrably affects parse compliance --- a reason for studies
collecting process descriptors \emph{by asking for them} to check, not a coefficient to apply. The
pilot also showed a \emph{third} effect in the P-T3 direction at the same cell (trace 9/10
on-prediction against bare $N = 21$ at 4/7, one-sided $p = 0.16$), so P-T3 replicated a pilot signal
rather than testing a blind prediction.

\subsection{Faithfulness with ground truth}

Claim and artifact sit in the same completion and the artifact is fully determined, so a method
line is checkable against the emitted coordinates. \texttt{arm\_f\_repro.trace\_faithfulness()}
matches numeric dimensions from the method line against the emitted layout signature, returning 38
of 41 scoreable claims matching (93\%) against the registered 90\% threshold. Nine of its 12
exclusions are regex coverage gaps, not claims without numeric content, and three are closer to
genuinely unscoreable (Appendix~E.8). The excluded set is enriched for non-canonical phrasings,
where a mismatch is a priori \emph{more} likely, so the exclusion was not conservative in a known
direction; worst case, had all 12 been unfaithful, the rate would have been $38/53 = 72\%$.

\textbf{Blind rescore under a frozen rubric.} The rubric (\texttt{p7\_faithfulness\_rubric.md}) was
committed to git as \texttt{e181d2a} at 03:56:56 on 2026-08-01, \textbf{before} any rescoring run,
so the freeze is checkable, and it pre-committed the reporting rule: below 90\%, P-T4 reported not
confirmed, full stop, no consolation framing. Scoring was blind --- claim text and circles only, no
predicted or rival values, no running tally, no threshold and no sight of the existing 38/41 result
--- with tallying afterward; all 60 labels are released in \texttt{p7\_blind\_labels.json}.
\textbf{Result:} over all 60 trace\_v2 rows, \textbf{54 MATCH, 2 MISMATCH, 4 UNSCOREABLE} --- a
corrected rate of \textbf{54/56 = 96.4\%} of scoreable claims, Wilson 95\% CI $[88\%, 99\%]$,
passed, and on 56 scoreable claims rather than 41. The CI's lower bound sits just below the
registered threshold, so ``clears 90\%'' is a statement about the point estimate at this $n$. The
two mismatches share a form, \emph{mis-descriptions of alternation}, and the original scorer's
three mismatches all score MATCH under hand adjudication (Appendix~E.11).

\textbf{Scope, three ways} (Appendix~E.11). The blind pass covers \emph{invalid} rows too, scoring
all 60 including the 7 invalid ones, 4 of the 60 unscoreable, where the original 93\% conditioned
on completions that produced a correct packing. It verifies the description against the
\emph{artifact}, not the process, so for a model already emitting grids it is closer to a
self-consistency check than to the causal question the chain-of-thought literature asks (\S7). And
it is confounded with the elicitation arm, computed where \S6.2 shows 87\% of valid outputs
concentrated onto one template and correct description is close to free, while the split by on-
versus off-prediction is not run here (\S9). What survives is narrow: in a domain where claims are
checkable, the method lines this proposer emits are, at 54 of 56 scoreable claims, true of the
object it produced.

